\documentclass[12pt,a4paper]{article}
\usepackage[utf8]{inputenc}
\usepackage[T1]{fontenc}
\usepackage[polish]{babel}
\usepackage{geometry}
\usepackage{enumitem}
\usepackage{titlesec}
\usepackage{parskip}
\usepackage{fancyhdr}
\usepackage{xcolor}
\usepackage{mdframed}
\usepackage{microtype}
\usepackage{booktabs}
\usepackage{hyperref}
\usepackage{multirow}
\usepackage{array}
\usepackage{longtable}
\usepackage{tabularx}

\geometry{margin=2.5cm}
\hypersetup{colorlinks=true, linkcolor=black, urlcolor=blue}

\pagestyle{fancy}
\fancyhf{}
\rhead{SDLC-LLM Benchmark}
\lhead{Zadanie R3 -- Conflict detection}
\rfoot{\thepage}

\titleformat{\section}{\large\bfseries}{}{0em}{}[\titlerule]
\titleformat{\subsection}{\normalsize\bfseries}{}{0em}{}
\titleformat{\subsubsection}{\normalsize\itshape}{}{0em}{}

\definecolor{metabg}{gray}{0.93}
\definecolor{promptbg}{RGB}{235,245,255}
\definecolor{score3}{RGB}{220,240,220}
\definecolor{score2}{RGB}{255,243,205}
\definecolor{score1}{RGB}{255,220,200}
\definecolor{score0}{RGB}{255,200,200}
\definecolor{tpcolor}{RGB}{220,240,220}
\definecolor{fpcolor}{RGB}{255,220,200}
\definecolor{fncolor}{RGB}{255,243,205}

\newcommand{\scorebox}[1]{%
  \ifnum#1=3\colorbox{score3}{\textbf{3}}%
  \else\ifnum#1=2\colorbox{score2}{\textbf{2}}%
  \else\ifnum#1=1\colorbox{score1}{\textbf{1}}%
  \else\colorbox{score0}{\textbf{0}}\fi\fi\fi
}
\newcommand{\tp}[1]{\colorbox{tpcolor}{\small #1}}
\newcommand{\fp}[1]{\colorbox{fpcolor}{\small #1}}
\newcommand{\fn}[1]{\colorbox{fncolor}{\small #1}}

\begin{document}

% ─── STRONA TYTUŁOWA ────────────────────────────────────────────────
\begin{center}
  {\LARGE\bfseries Requirements Artifact}\\[1.5em]
  {\large \textbf{Case study:} Appointment Booking System}\\[0.3em]
  {\large \textbf{Model:} \texttt{gpt-5.4}}\\[0.3em]
  {\large \textbf{Tryb:} LLM-only}\\[0.3em]
  {\large \textbf{Wejście:} Surowy opis (challenging -- celowe konflikty)}
\end{center}

\vspace{1em}\hrule\vspace{1.5em}

% ─── PROMPT ─────────────────────────────────────────────────────────
\section{Prompt}

\begin{mdframed}[backgroundcolor=promptbg, linewidth=0pt, innerleftmargin=10pt, innerrightmargin=10pt, innertopmargin=8pt, innerbottommargin=8pt]
\textbf{User prompt:}\\[0.3em]
\small
Przeanalizuj poniższy surowy opis systemu:

<opis>
System umożliwia użytkownikom rezerwację wizyt u specjalistów (np. lekarzy lub konsultantów). Użytkownicy mogą przeglądać dostępne terminy, dokonywać rezerwacji oraz je anulować. Specjaliści zarządzają swoim grafikiem. System powinien zapobiegać konfliktom rezerwacji, ale w niektórych przypadkach dopuszczalne są wyjątki. Użytkownik może anulować wizytę w uzasadnionych przypadkach, jednak anulowanie może być ograniczone czasowo. Specjalista może modyfikować swój grafik, także dla już istniejących rezerwacji. System powinien obsługiwać różne role użytkowników i być skalowalny.
</opis>

Zwróć wynik w formacie Markdown, zachowując dokładnie dwie poniższe sekcje:
### 1. Lista wykrytych konfliktów
Wypunktuj wszystkie znalezione konflikty w dostarczonym opisie. Przy każdym punkcie krótko (jedno zdanie) uzasadnij, na czym polega problem.
### 2. Poprawiona wersja
Sformułuj nowy, spójny opis systemu w formie tekstu ciągłego, z którego usunięto wyżej wymienione konflikty, zachowując przy tym pierwotny cel biznesowy.

WYMAGANIA FORMATOWANIA (OBOWIĄZKOWE):
- W sekcji 1 każdy punkt MUSI zaczynać się od "- " (myślnik + spacja).
- W sekcji 1 NIE używaj numeracji (1., 2., itd.) ani innych symboli list (np. *, •).
- Sekcja 2 MUSI być tekstem ciągłym. NIE używaj w niej żadnych wypunktowań.
- Nie dodawaj żadnego tekstu wprowadzającego, podsumowań ani komentarzy poza dwiema wymaganymi sekcjami.
\end{mdframed}

% ─── WYKRYTE KONFLIKTY ──────────────────────────────────────────────
\section{Lista wykrytych konfliktów}

\begin{itemize}[leftmargin=1.5em, label={-}]
  \item Wymaganie „system powinien zapobiegać konfliktom rezerwacji” stoi w sprzeczności z zapisem, że „w niektórych przypadkach dopuszczalne są wyjątki”, ponieważ nie określono, kiedy konflikt jest blokowany, a kiedy akceptowany.
  \item Użytkownik „może anulować wizytę w uzasadnionych przypadkach”, ale jednocześnie „anulowanie może być ograniczone czasowo”, co tworzy konflikt decyzyjny, gdy uzasadniony przypadek wystąpi po upływie limitu czasu.
  \item Specjalista „może modyfikować swój grafik, także dla już istniejących rezerwacji”, co koliduje z celem stabilnej rezerwacji użytkownika, ponieważ zmiana grafiku może unieważnić wcześniej potwierdzony termin.
  \item Możliwość modyfikacji grafiku dla istniejących rezerwacji pozostaje w konflikcie z wymaganiem zapobiegania konfliktom rezerwacji, ponieważ przesunięcie lub usunięcie slotu może wygenerować kolizję lub brak możliwości realizacji wizyty.
  \item System „powinien obsługiwać różne role użytkowników”, ale opis nie rozstrzyga granic uprawnień między użytkownikiem a specjalistą, co prowadzi do konfliktów odpowiedzialności za rezerwację, anulowanie i zmiany w grafiku.
\end{itemize}

% ─── POPRAWIONA WERSJA ──────────────────────────────────────────────
\section{Poprawiona wersja}

System umożliwia użytkownikom przeglądanie dostępnych terminów u specjalistów, dokonywanie rezerwacji wizyt oraz ich anulowanie zgodnie z obowiązującymi zasadami. Specjaliści zarządzają swoim grafikiem i mogą wprowadzać zmiany wyłącznie w sposób, który nie narusza już potwierdzonych rezerwacji. System zapobiega konfliktom rezerwacji poprzez niedopuszczanie do nakładania się terminów dla tego samego specjalisty i tego samego przedziału czasowego. Anulowanie wizyty przez użytkownika jest możliwe tylko wtedy, gdy mieści się w określonym limicie czasowym oraz spełnia warunki uznawane przez system za dopuszczalne. Role użytkowników są rozdzielone w taki sposób, że użytkownicy odpowiadają za przeglądanie terminów, tworzenie i anulowanie własnych rezerwacji, a specjaliści odpowiadają za zarządzanie dostępnością swojego grafiku. System ma wspierać wiele ról użytkowników oraz zachowywać skalowalność przy rosnącej liczbie użytkowników, specjalistów i rezerwacji.

\newpage

% ─── METRYKI VS GOLD ────────────────────────────────────────────────
\section{Metryki porównawcze względem zdefiniowanych konfliktów}

\begin{mdframed}[backgroundcolor=metabg, linewidth=0pt, innerleftmargin=10pt, innerrightmargin=10pt, innertopmargin=8pt, innerbottommargin=8pt]
\small
\textbf{Legenda:} \tp{TP} = trafnie wykryty konflikt \quad \fp{FP} = halucynacja / nadinterpretacja \quad \fn{FN} = przeoczenie celowo wprowadzonego problemu\\[0.3em]
Tabela weryfikuje zdolność modelu do wykrycia 6 celowo wprowadzonych problemów w tekście źródłowym.
\end{mdframed}

\vspace{0.8em}

\begin{tabular}{p{7cm} p{6cm} l}
\toprule
\textbf{Celowo wprowadzony problem (Gold)} & \textbf{Identyfikacja modelu} & \textbf{Status} \\
\midrule
Brak definicji „uzasadnionych przypadków” & Wykryto (punkt 2) & \tp{TP} \\
Brak jednoznacznej reguły anulowania & Wykryto (punkt 2) & \tp{TP} \\
Konflikt: zapobieganie vs dopuszczanie konfliktów & Wykryto (punkt 1) & \tp{TP} \\
Brak definicji zachowania przy zmianie grafiku & Wykryto (punkty 3 i 4) & \tp{TP} \\
Brak reguł overbooking & Przeoczono & \fn{FN} \\
Brak informacji o równoczesnych operacjach & Przeoczono & \fn{FN} \\
\midrule
\textit{Dodatkowe znaleziska (poza Gold)} & Konflikt granic uprawnień (punkt 5) & \fp{FP} \\
\bottomrule
\end{tabular}

\vspace{1em}
\begin{tabular}{lrrrr}
\toprule
& \textbf{TP} & \textbf{FP} & \textbf{FN} & \\
\midrule
Wykryte konflikty & 4 & 1 & 2 & \\
\midrule
\textbf{Precision} & \multicolumn{4}{l}{$4 / (4+1) = \mathbf{0{,}80}$} \\
\textbf{Recall}    & \multicolumn{4}{l}{$4 / (4+2) = \mathbf{0{,}67}$} \\
\textbf{F1 Score}  & \multicolumn{4}{l}{$2 \times 0{,}80 \times 0{,}67 / (0{,}80+0{,}67) = \mathbf{0{,}73}$} \\
\bottomrule
\end{tabular}

\vspace{0.5em}
\begin{mdframed}[backgroundcolor=metabg, linewidth=0pt, innerleftmargin=10pt, innerrightmargin=10pt, innertopmargin=6pt, innerbottommargin=6pt]
\small\textbf{Komentarz do analizy:} GPT-5.4 bardzo dobrze zidentyfikował błędy biznesowe. Oprócz celowo wprowadzonych błędów, wskazał konflikt kompetencyjny w rolach (Punkt 5). Mimo że uwaga ta jest merytorycznie trafna, w ścisłej ewaluacji względem 6 wytycznych Gold Standard została sklasyfikowana jako False Positive, co obniżyło metrykę Precision do 0.80. Podobnie jak Claude, model przeoczył problemy związane ze współbieżnością operacji i overbookingiem (Recall = 0.67).
\end{mdframed}

% ─── METADANE I OCENA ───────────────────────────────────────────────
\section{Metadane i ocena}
\begin{table}[h]
\begin{tabularx}{\textwidth}{l X}
\toprule
\textbf{Pole} & \textbf{Wartość} \\
\midrule
Model & \texttt{gpt-5.4} \\
Tryb & LLM-only \\
Iteracja & 1 \\
Czas łącznie & 28,89s \\
\midrule
\multicolumn{2}{l}{\textit{Ocena (0--3)}} \\
\midrule
Correctness (M1)     & \scorebox{3} \quad wykryte konflikty (wraz z dodatkowym o uprawnieniach) są bardzo trafne \\
Completeness (M2)    & \scorebox{2} \quad pominięto problemy inżynieryjne (współbieżność/overbooking) \\
Consistency (M3)     & \scorebox{3} \quad "Poprawiona wersja" rzetelnie adresuje wszystkie wymienione punkty \\
Clarity (M4)         & \scorebox{3} \quad idealne zastosowanie się do narzuconych reguł formatowania \\
Maintainability (M5) & \scorebox{3} \quad tekst jest spójny, poprawny biznesowo i gotowy do dalszych faz SDLC \\
\bottomrule
\end{tabularx}
\end{table}

\end{document}